% Capítulo de la tesis (texto derivado de envios/paper_empirico/cuerpo.tex; se edita aquí).
\chapter{Conclusiones}
\label{sec:conclusiones}

Esta tesis preguntó si el mercado penaliza de forma no lineal la coexistencia de fragilidad bancaria y vulnerabilidad macroeconómica, de modo que la coincidencia de una cola adversa del crecimiento y una alta fragilidad bancaria amplifique el spread soberano de las economías emergentes. Para responderla articuló en un solo marco dos literaturas que se han mantenido separadas ---la del nexo banca--soberano y la del riesgo a la baja del crecimiento--- y construyó, para trece economías emergentes, las dos métricas que lo operacionalizan: la fragilidad bancaria sistémica a partir de 113 bancos extraídos homogéneamente de Bloomberg, y el riesgo de cola del crecimiento por regresión cuantílica en pseudo--tiempo real.

La respuesta tiene tres partes. Primero, ambos riesgos se asocian con un spread mayor: la elasticidad del spread a la fragilidad bancaria es de 0,10 ---unos 4,7 puntos básicos por unidad de $JLoss$--- y un punto porcentual adicional de riesgo de cola se asocia con un spread un 3\% mayor; la asociación de la fragilidad replica la que documenta la literatura más próxima. Segundo, esas asociaciones no resisten como efectos causales: el spread retroalimenta a la fragilidad, ambos son muy persistentes, y con la inferencia que exige un panel de trece países ninguna es significativa. Tercero, en el promedio del panel la complementariedad entre ambos riesgos no va más allá de la que implica una estructura proporcional; la amplificación adicional aparece solo en momentos y economías específicas ---fuera de las crisis con respaldo oficial y en las economías cuya deuda fija un inversionista extranjero---, con evidencia que resiste la inferencia con pocos \textit{clusters} pero no los controles en logaritmos ni el hecho de que ese subconjunto se definiera \textit{ex post}.

La hipótesis con la que se aprobó esta investigación queda, así, sin respaldo en su versión universal ---una simultaneidad que se manifieste siempre y en todas las economías--- y con respaldo condicional en una versión acotada: la amplificación sería un fenómeno de ciertos regímenes y de cierto tipo de economía. Esa distinción es también la implicancia de política más relevante: más que una regulación uniforme para todas las economías en todo momento, la evidencia apunta a caracterizar los casos en que la doble vulnerabilidad se traslada al precio de la deuda soberana, a condición de que esa evidencia se confirme con más datos.

El aporte del trabajo es delimitar con precisión qué permite afirmar este diseño y qué no, mostrar dónde buscar la amplificación y por qué el panel disponible no basta para establecerla, y dejar construida la infraestructura de datos ---reproducible y trazable, en \texttt{1\_Codigo/Panel/bbg/}--- sobre la que un diseño con más economías, más frecuencia o variación exógena podría hacerlo.
